\chapter{Resultados}
\label{chap6}

Atualmente, o sistema encontra-se em fase de homologação. Essa etapa tem como objetivo validar, junto aos usuários envolvidos no processo de gestão do PAA e ao Assentamento PDS Osvaldo de Oliveira, se as funcionalidades implementadas atendem às necessidades identificadas durante o levantamento de requisitos e se o fluxo do sistema corresponde à realidade operacional do programa.

A versão em homologação contempla as principais funcionalidades necessárias para a criação e o acompanhamento de um edital, permitindo organizar informações que antes eram controladas de forma manual, principalmente por meio de planilhas. Entre essas funcionalidades, destacam-se o cadastro de usuários, a aprovação de acesso ao sistema, o cadastro de famílias, a criação de editais e a definição dos perfis de acesso dos participantes.

Com isso, é possível observar que o sistema propõe uma nova forma de olhar para a gestão do PAA. Ao centralizar as informações em uma plataforma web, o acompanhamento dos editais tende a se tornar mais objetivo, claro e transparente. Além disso, a ferramenta reduz a dependência de controles manuais e facilita a consulta das informações pelos diferentes usuários envolvidos no processo.

Dessa forma, mesmo ainda estando em homologação, o sistema já demonstra potencial para contribuir com a organização da gestão do PAA no assentamento. A validação com os usuários será fundamental para identificar ajustes necessários, melhorar a experiência de uso e garantir que a ferramenta esteja adequada às demandas reais da comunidade antes de sua implantação definitiva.


